\paragraph{Independent normalization of the complete-line spectral density.}
The Laplacian is positive, its spectral parameter is
$\lambda=1/4+t^2$, and all densities below are with respect to
$dt$ on $t>0$, not with respect to $d\lambda$.
In Fermi coordinates about a complete geodesic the metric is
$dr^2+\cosh^2r\,dv^2$.  Integrating the hyperbolic resolvent
kernel over that geodesic gives the transverse resolvent for
\[
 L=-\partial_r^2-\tanh r\,\partial_r
 \quad\hbox{on }L^2(\mathbb R,\cosh r\,dr)
\]
with source $\delta_0$; the area-density factor at the source
is $\cosh0=1$.  For real $s>1$ the integrated resolvent is
absolutely convergent, and uniqueness of the decaying resolvent
solution identifies it with the expression below.  Analytic
continuation then gives its outgoing and incoming boundary values.

Multiplication by $\sqrt{\cosh r}$ is unitary to $L^2(dr)$ and
conjugates $L-1/4$ to
\[
 H=-\partial_r^2+V(r),\qquad
 V(r)=\tfrac14\operatorname{sech}^2r.
\]
Retain the normalization
\[
 f_s(r)=2^{-s}\cosh(r)^{-s}
 {}_2F_1\!\left(\frac s2,\frac s2;s+\frac12;\operatorname{sech}^2r\right),
 \qquad f_s(r)\sim e^{-sr}.
\]
Since the potential is even, the Green kernel with one argument at
zero is
\[
 G_s(r,0)=-\frac{f_s(|r|)}{2f_s'(0)}.
\]
Indeed it is continuous and its first derivative has jump $-1$;
thus $(L-s(1-s))G_s=\delta_0$ with the required sign.  Equivalently,
if $j_+(r,t)=\sqrt{2\cosh r}\,f_{1/2-it}(r)\sim e^{itr}$
and $j_-(r,t)=j_+(-r,t)$, their Wronskian at zero is
$2j_+(0,t)j_+'(0,t)$; the Green kernel for $H-t^2$ is
$-j_+(r,t)/(2j_+'(0,t))$ for $r>0$.

The upper spectral boundary corresponds to $s=1/2-it$:
$s=1/2+\epsilon-it$ gives
$s(1-s)=1/4+t^2-\epsilon^2+2i\epsilon t$.
Stone's resolvent-jump formula and $d\lambda=2t\,dt$ therefore give
\[
 \boxed{\quad
 \mathcal P_t(r)=\frac{t}{\pi i}
 \left(G_{1/2-it}(r,0)-G_{1/2+it}(r,0)\right).
 \quad}
\]
There is no additional factor of two from the two sides of the
geodesic; both sides were already included in the derivative jump.
The denominators do not vanish for $t>0$, since otherwise both
conjugate derivatives would vanish at zero, contradicting the
nonzero incoming--outgoing Wronskian at infinity.

For an exact formula and an independent sign check, DLMF (15.8.4),
whose principal-branch hypotheses hold for
$z=\operatorname{sech}^2r\in(0,1)$, gives the coefficient of
$(1-z)^{1/2}$ at $z=1$.  Since
$(1-z)^{1/2}=\tanh r$ for $r>0$,
\[
 f_s'(0)=-\frac{2^{1-s}\sqrt\pi\,\Gamma(s+1/2)}{\Gamma(s/2)^2}.
\]
It follows, with the previously defined $K_{s,0}$ and $k_s$, that
the exact sharp-source kernel can also be written
\[
 q_t(u)=\frac{K_{s,0}}2 k_s(u)
       =\frac{a}{2\cosh\tau(u)}
          \frac{f_s'(\tau(u))}{f_s'(0)}
       =-\frac{a}{\cosh\tau(u)}G_s'(\tau(u),0),
 \qquad s=\tfrac12+it.
\]
All gamma factors cancel exactly in this identity.

\paragraph{The leading coefficient and uniform symbol.}
For $r\geq0$, define the normalized outgoing amplitude
\[
 \alpha_t(r)=
 \frac{it\,e^{-itr}\sqrt{\cosh r}\,f_{1/2-it}(r)}
      {f_{1/2-it}'(0)}.
\]
The exact spectral-density formula becomes
\[
 \mathcal P_t(r)=\frac{
    \alpha_t(r)e^{itr}+\overline{\alpha_t(r)}e^{-itr}}
                  {2\pi\sqrt{\cosh r}}.
\]
The uniform Volterra estimates, now applied on the entire half-line,
show
\[
 \alpha_t(r)=1+O(t^{-1}),\qquad
 \partial_r\alpha_t(r)=O(e^{-2r}),
 \qquad t\geq1, r\geq0.
\]
For explicit control, put
$J(r)=\int_r^\infty V(v)\,dv=(1-\tanh r)/4$.
For the unit-normalized Jost solution $j_+$,
\[
 |j_+(r,t)-e^{itr}|\leq e^{J(r)/t}-1,
 \qquad
 |j_+'(r,t)-it e^{itr}|\leq J(r)e^{J(r)/t}.
\]
In particular
\[
 |j_+'(0,t)|\geq t-\tfrac14e^{1/4}
 \geq\left(1-\tfrac14e^{1/4}\right)t>0
 \qquad(t\geq1).
\]
These estimates give the displayed bounds for $\alpha_t$ directly,
including its derivative after the outgoing exponential is removed.
Successive Volterra expansions give the full expansion to every
fixed inverse power of $t$, with a remainder uniform in $r\geq0$.
The nonconstant coefficient derivatives decrease exponentially,
as required for the Lipschitz factors in the endpoint BV lemma.
Consequently the leading term is exactly
\[
 \boxed{\quad
 \mathcal P_t(r)=\frac{\cos(tr)}{\pi\sqrt{\cosh r}}
       +O\!\left(\frac{1}{t\sqrt{\cosh r}}\right).
 \quad}
\]
There is no fixed scattering phase or sign which could cancel the
stationary-phase coefficient.  The free one-dimensional case
$V=0$ gives $\pi^{-1}\cos(tr)$, checking the $dt$ normalization.

There is also a positive exact check at the source.  The connection
formula gives
\[
 G_s(0,0)=\frac14
 \left(\frac{\Gamma(s/2)}{\Gamma((s+1)/2)}\right)^2.
\]
Gamma reflection, with $z=1/4-it/2$, yields
\[
 \mathcal P_t(0)=
 \frac{t\sinh(\pi t)}{4\pi^3}
       |\Gamma(1/4+it/2)|^4>0.
\]
Stirling's formula makes this tend to $1/\pi$, consistently with
the full-line normalization and the preceding uniform estimate.

\paragraph{Source scope.}
DLMF equation (15.8.4) and its branch conventions were opened at
\texttt{https://dlmf.nist.gov/15.8} on 15 November 2026.
The resolvent derivative jump, unitary conjugation and
spectral-variable Jacobian have been derived above rather than
imported with an unspecified spectral-measure convention.
